\section{Experimental Methodology}
\label{sec:experimental-methodology}

\subsection{Research Questions}
\label{sec:rqs}

\begin{itemize}
  \item \textbf{RQ1 (quality).} On ISD-Agent-Bench, how does AlignmentGraph-ISD
  compare with representative ISD agents in artifact quality, under both a large
  model and a small language model?
  \item \textbf{RQ2 (ablation).} Which harness mechanism drives that quality, and
  does its contribution change with model scale?
\end{itemize}

\subsection{Dataset}
\label{sec:dataset}

We use ISD-Agent-Bench~\cite{isdagentbench}, whose scenarios span a
Context~$\times$~ISD matrix. Evaluation is run on the \texttt{test\_90} split---90
of the test scenarios, stratified over 10 domains $\times$ 3 difficulty levels
$\times$ 3 scenarios per cell (seed 42)---which gives broad, balanced coverage at a
tractable judging cost.

\subsection{Models and Setup}
\label{sec:setup}

To probe model independence we run every configuration at two distinct points on
the model-scale spectrum: \textbf{DeepSeek-v4-flash} (a large LLM) and
\textbf{Qwen3.5-9B} (a 9B SLM), with identical configurations apart from the agent
model. Agent and judge models are configured separately to avoid self-preference
bias; the judges are held fixed across all runs and use the benchmark's default
multi-judge configuration (\texttt{openai/gpt-4o-mini} and
\texttt{google/gemini-2.5-flash-lite}). Each run records a manifest (model, judges,
prompt versions, commit, seed, decoding configuration) for reproducibility.

\subsection{Conditions}
\label{sec:conditions}

For RQ1, AlignmentGraph-ISD is integrated as a seventh agent under the benchmark's
input/output schema and compared with six agents: Baseline (single prompt),
ReAct-ADDIE, ADDIE-Agent, Dick-Carey-Agent, RPISD-Agent, and EduPlanner.

For RQ2, the ablation removes exactly one mechanism at a time from the full
harness. Each arm is a self-contained workflow (no shared configuration flags), so
no mechanism can leak between arms:
\begin{itemize}
  \item \texttt{full} --- multi-agent decomposition + per-agent verify/repair;
  \item \texttt{wo-ma} --- decomposition removed: a single agent generates the
  whole program and repairs itself holistically (verify/repair retained);
  \item \texttt{wo-qc} --- verify/repair removed: the same fixed-order specialists
  generate, with no quality-control tools.
\end{itemize}
Comparing \texttt{full} against \texttt{wo-ma} isolates multi-agent decomposition;
against \texttt{wo-qc} isolates per-agent verification/repair.

\subsection{Metrics}
\label{sec:metrics}

The primary metrics are the benchmark's: a 100-point ADDIE artifact rubric and a
trajectory score, combined into a composite \emph{Total}, each a two-judge mean.
We additionally report a deterministic \emph{residual alignment-defect audit}: the
verifier is replayed over each arm's final output and the surviving content-level
violations are counted. This measures problem condition~(i) directly, capturing a
property the rubric does not explicitly reward.

\subsection{Statistical Analysis}
\label{sec:stats}

Comparisons are paired by scenario. We use the Wilcoxon signed-rank test
(tie-corrected normal approximation) on the composite score, with Holm correction
across the comparisons in each run, and report the mean paired difference with a
bootstrap 95\% confidence interval (10{,}000 resamples). For RQ2 the contribution
of a mechanism is $\texttt{full} - \texttt{arm}$, paired by scenario.

\subsection{Reproducibility}
\label{sec:reproducibility}

ISD-Agent-Bench and its six baselines are public. Every run is pinned by a manifest
(\cref{sec:setup}); the three ablation arms are separate workflow modules selected
by name; and the analysis scripts (paired tests and the residual audit) are
released with the paper.
